% BrainAI - built by Mohammed Zahid.
% Compiles with pdflatex on any standard TeX distribution, or on Overleaf as-is.
\documentclass[10pt,twocolumn]{article}

\usepackage[margin=0.75in]{geometry}
\usepackage{amsmath,amssymb}
\usepackage{booktabs}
\usepackage{graphicx}
\usepackage{hyperref}
\usepackage{xcolor}
\usepackage{caption}
\usepackage{listings}
\usepackage{microtype}

\hypersetup{colorlinks=true, linkcolor=blue!50!black,
            citecolor=blue!50!black, urlcolor=blue!50!black}

\captionsetup{font=small,labelfont=bf}

\lstset{basicstyle=\ttfamily\scriptsize, breaklines=true,
        frame=single, framesep=3pt, columns=fullflexible}

\newcommand{\win}{\textbf{$\checkmark$}}
\newcommand{\lose}{$\times$}

\title{\vspace{-1.5em}\textbf{BrainAI: A Spiking Cognitive Architecture for\\
Embodied Continual Learning}\\[0.35em]
\large Every subsystem measured against a control, including the two that lost}

\author{
  Mohammed Zahid Imityaz Wadiwale\\
  \small \texttt{zahid.wadiwale1234@gmail.com}
}
\date{}

\begin{document}
\maketitle

\begin{abstract}
\noindent
We present \textbf{BrainAI}, a spiking multimodal cognitive architecture for
embodied continual learning. A single model perceives text, images, video,
audio and a typed sensor suite; routes each input through a reflex arc, an
always-on subconscious pathway and a gated conscious pathway; maintains five
distinct memory systems split along the differentiability line; predicts future
latent states and learns from prediction error without labels; weighs new
claims against evidence before believing them; pursues interruptible goals;
plans by imagination; updates its own weights during inference under bounded
plasticity; and consolidates offline during a sleep cycle. Motor commands pass
through a non-learned safety controller that the learning system provably
cannot modify.

The methodological commitment of this work is that \emph{every} subsystem is
measured against the obvious simpler alternative, and every result is reported.
Measured outcomes include a $2.16\times$ inference speedup from genuine
computational gating, a $74\%$ reduction in catastrophic forgetting, $3.7\times$
less degradation over 400 inference-time weight updates, a belief system that
defeats a naive vote-counter on four of five adversarial scenarios, and
collusion detection that holds a protected axiom at confidence $1.0000$ against
a 600-identity Sybil attack that otherwise drives it to $0.0001$.

The two proposals with mixed or negative results are reported in full. The
architecture's own activation function, AMSA, is outperformed by SwiGLU on
perplexity, wall-clock, and the continual-learning problem it was designed for,
with the loss gap \emph{widening} with scale---but on the axis it was actually
built for, efficiency, it \emph{wins}: at comparable loss it fires fewer spikes
and activates $\sim$14 points less of the feed-forward layer, across three seeds,
making it a quality--efficiency tradeoff rather than a flat loss. The world model
is outperformed at one-step prediction by a linear map, though the deficit does
not propagate to multi-step rollout or planning, where the two tie, and only the
learned model reports calibrated uncertainty ($r=0.565$).

We additionally catalogue fourteen \emph{silent} failure modes discovered only
because controls were run---including an elastic-weight-consolidation penalty
that contributed $10^{-4}$ against a task loss of $2.5$ and therefore did
nothing while appearing to work, and a rate limiter disabled on every call by a
tensor-shape comparison. We argue this catalogue is among the more transferable
contributions of the work.
\end{abstract}

\section{Introduction}

Contemporary large models separate training from deployment: weights are fixed
at inference, knowledge is frozen at a cutoff, and adaptation happens through
context or retrieval rather than through the network itself. For a stationary
task distribution this is the right engineering trade. For an embodied agent
expected to operate for years in one environment it is not: such an agent must
acquire knowledge after deployment, retain what it already knows, decide which
of the things it is told are actually true, and do all of this without a
training loop supervising it.

BrainAI is an attempt to build that system end to end. It is not a language
model with memory attached; it is a cognitive architecture in which the
language head is one of several outputs. The design draws on complementary
learning systems theory~\cite{mcclelland1995}, dual-process
accounts of cognition~\cite{kahneman2011}, and the spiking-network
literature~\cite{maass1997,neftci2019}, and assembles components---spiking
transformers, mixture-of-experts routing, elastic weight
consolidation~\cite{kirkpatrick2017}, experience replay~\cite{rolnick2019},
latent dynamics models~\cite{hafner2019}---that exist individually but, to our
knowledge, have not previously been integrated with the specific interfaces
described here.

\subsection{The problem with architecture papers}

An architecture paper is easy to write badly. A diagram with a dozen boxes, a
plausible mechanism for each, and a benchmark on the one task where the whole
thing happens to work is a well-established genre, and it is nearly
unfalsifiable. The failure mode is not dishonesty; it is that a component which
does nothing at all looks exactly like a component that works, unless something
is run without it.

We therefore adopted a single rule: \textbf{no component is described without a
control}. For each subsystem we identified the obvious simpler alternative---the
thing a reasonable engineer would do instead---implemented it, and measured
both. This rule is the reason this paper contains negative results, and it is
also the reason we trust the positive ones.

\subsection{Contributions}

\begin{enumerate}\itemsep2pt
  \item \textbf{An integrated architecture} (\S\ref{sec:arch}) whose novelty is
        in the interfaces between components rather than in any single one: a
        memory hierarchy split by differentiability, a prediction ledger that
        grounds epistemic trust in physical observation, and a safety layer
        structurally unreachable by the learning system.
  \item \textbf{A controlled evaluation} comprising thirteen comparisons
        (\S\ref{sec:results}), each against a simpler alternative.
  \item \textbf{Two negative results} reported in full
        (\S\ref{sec:negative}), including a scale sweep testing---and
        rejecting---the standard defence that the losing component merely needs
        a larger model.
  \item \textbf{A catalogue of fourteen silent failure modes}
        (\S\ref{sec:silent}) found only by running controls. Several are
        general to continual-learning and robotics implementations rather than
        specific to this one.
\end{enumerate}

All results are produced by scripts in the accompanying repository. Every
number in this paper, including those in \S\ref{sec:negative}, is reproducible
on a CPU.

\section{Related Work}

\paragraph{Continual learning.} Catastrophic interference has been understood
since~\cite{mccloskey1989}. Regularisation approaches penalise movement in
parameters important to earlier tasks~\cite{kirkpatrick2017,zenke2017};
rehearsal approaches interleave stored examples~\cite{rolnick2019}. We use both,
and find (\S\ref{sec:continual}) that the regularisation term is worthless
unless its scale is calibrated---a detail rarely reported.

\paragraph{Spiking transformers.} Surrogate-gradient training~\cite{neftci2019}
made deep spiking networks practical; adaptive-threshold neurons~\cite{bellec2018}
supply temporal memory without recurrence. Spiking transformer variants have
been demonstrated~\cite{zhou2023}. Our contribution here is not the block
design but two implementation constraints we found to be load-bearing
(\S\ref{sec:silent}).

\paragraph{World models.} PlaNet~\cite{hafner2019} and Dreamer~\cite{hafner2020}
established latent dynamics models for planning. Prediction error as intrinsic
motivation dates to~\cite{pathak2017}; the noisy-TV pathology and its
diagnosis are from~\cite{burda2019}. We compare against both an RSSM-family
recurrent baseline and---more informatively---a linear one.

\paragraph{Epistemics.} Sybil attacks on reputation systems are
classical~\cite{douceur2002}. To our knowledge the specific combination used
here---log-odds pooling with per-source diminishing returns, credibility
shrinkage for unvetted sources, and behavioural collusion clustering---has not
been applied to a robot's belief store.

\section{Architecture}
\label{sec:arch}

\subsection{Spiking substrate}

Spikes are true Heaviside steps; the forward pass is binary, so reported spike
rates are genuine sparsity rather than a soft approximation. Training uses an
arctan surrogate on the backward pass only. Neurons are adaptive
leaky-integrate-and-fire: each spike raises that neuron's own threshold, decaying
over $\tau_{\text{adapt}}$, which supplies temporal memory across simulation
timesteps without recurrent weights.

Two constraints proved necessary. First, adaptation is
$a[t] = \rho\,a[t{-}1] + s[t{-}1]$ with $a[0]=s[0]=0$, so at $T=2$ the decay
$\rho$ only ever multiplies zero and threshold adaptation is inert; $T\geq3$ is
the minimum useful configuration. Second, the residual stream must remain
analogue---making it spiking destroys the gradient path and the model does not
train.

\subsection{Latency-tiered pathways}

Three pathways of increasing cost run in parallel. The \emph{reflex} arc is one
shallow stack from perception to a candidate output; it never consults memory.
The \emph{subconscious} pathway always runs and emits a salience score. The
\emph{conscious} pathway is wide and expert-routed~\cite{shazeer2017,fedus2022}
and runs only when salience clears a threshold.

Blocks are pre-norm transformer blocks~\cite{vaswani2017} with binary spike
inputs to the attention projections. Attention offers softmax or
sparsemax~\cite{martins2016}, the latter producing exactly-zero weights so the
attention pattern is sparse as well.

The gate is soft during training, so the salience head receives gradient and
learns when deliberation pays, and hard at inference, where batch rows below
threshold are selected out and the wide stack genuinely does not execute. This
distinction matters: a gate that computes everything and multiplies by zero is
observationally identical from outside, which is why it is measured by
wall-clock (\S\ref{sec:results}).

\subsection{Memory hierarchy}

Five memories, split along the differentiability line
(Table~\ref{tab:memory}). The split is the anti-forgetting strategy: knowledge
acquired after training enters the non-differentiable \emph{stores}, where
writing it cannot damage a weight an older skill depends on. Only slow,
importance-gated consolidation moves anything from a store into the weights.

\begin{table}[t]
\centering\small
\begin{tabular}{@{}llc@{}}
\toprule
Memory & Backing & Differentiable \\
\midrule
Sensory & deque & no \\
Short-term & spiking LSTM & yes \\
Working & attention over slots & yes \\
Episodic & vector store & no \\
Semantic & vector store + graph & no \\
\bottomrule
\end{tabular}
\caption{The memory hierarchy. Neural state for reasoning; addressable stores
for facts.}
\label{tab:memory}
\end{table}

Two design decisions are load-bearing. The key and value storage projections
are \textbf{frozen random projections, not learned}: every key already in the
store was written with those weights, so training them would silently
invalidate the entire store as old keys and new queries drifted out of a shared
space. A random projection preserves the cosine geometry retrieval requires.
Second, retrieval is \textbf{non-differentiable by design}, so a bad memory can
mislead the model but can never corrupt a gradient. The cost is that the
importance scorer receives no signal from the language loss; it is instead
trained during sleep against a prioritised-replay target.

\subsection{World model}

Given a latent state and candidate action, the model predicts the next latent
state and, separately, the magnitude of its own error. Three capabilities
derive from that one mechanism:

\begin{itemize}\itemsep1pt
  \item \textbf{Grounding.} A \emph{prediction ledger} lets a belief stake a
        claim on a future observation. When the observation arrives, the ledger
        resolves it and reports which sources were vindicated.
  \item \textbf{Curiosity.} Prediction error normalised against \emph{expected}
        error. The second head is what distinguishes ``I do not know'' from
        ``this is inherently random'', and is why the drive does not park in
        front of a noise source~\cite{burda2019}.
  \item \textbf{Planning.} Rollouts occur in imagination; rollout confidence
        compounds multiplicatively, so a long plan through poorly understood
        dynamics is correctly distrusted.
\end{itemize}

\subsection{The Reality Cortex}

Salience and truth are different questions. A claim repeated by a thousand
people scores highly on the first and may be false. The Reality Cortex sits on
the knowledge path---episodic $\rightarrow$ validate $\rightarrow$ semantic
$\rightarrow$ sleep $\rightarrow$ weights---and not in the token path, since
contradiction detection over a belief graph is the wrong operation to run per
token.

Evidence is pooled in log-odds. Each observation contributes
$\log\!\big(t/(1-t)\big)$ for its source's reliability $t$, with refutation the
same quantity negated. Three discounts then apply:

\begin{align}
\text{within-source:} \quad & w_n = 1/n \\
\text{across-source:} \quad & w_k = 1/\sqrt{k} \\
\text{credibility:} \quad & c = \tfrac{n_{\text{eff}}}{n_{\text{eff}} + \kappa}
\end{align}

The first makes repetition from one voice saturate harmonically. The second
treats $k$ agreeing but unvetted sources as fewer than $k$ independent
witnesses, since correlated evidence pooled as independent overcounts by
roughly $\sqrt{k}$. The third shrinks a source's likelihood ratio until it has
a track record, so a freshly minted identity does not arrive at full strength.

For functional relations---those admitting one object, like
\texttt{capital\_of}---confidence is rival-relative,
$\sigma(\ell_i - \max_{j\neq i}\ell_j)$, so two equally supported alternatives
sit at $0.5$ and neither is promoted.

\paragraph{Collusion detection.} The above still assumes distinct names are
distinct witnesses, which is precisely what a Sybil attack
violates~\cite{douceur2002}. A puppet ring is visible in its behaviour: the
accounts appear on the same claims with the same polarity. Sources whose voting
records match are clustered and cast one vote between them, detected by exact
signature grouping ($O(S)$) followed by Jaccard refinement on a bounded sample.

\subsection{Executive, planner, uncertainty, safety}

The \textbf{executive} maintains a goal stack scored by
$\text{priority}\times\text{urgency}\times\text{feasibility}$, handles
interrupts that save and restore working-memory contents, suspends conflicting
goals rather than failing them, and schedules sleep from four independent
pressures with idleness as a multiplier rather than a term.

The \textbf{planner} is HTN-style~\cite{erol1994}: a method library maps goals
to ordered subgoals recursively. Candidate plans are scored by world-model
rollout. Procedures that survive the Reality Cortex are imported as executable
methods, which is where evidence and capability meet.

\textbf{Uncertainty} is carried per stage and composed multiplicatively with a
floor---a chain is no more trustworthy than its weakest link. Sources include
predictive entropy, spike-rate deviation from a layer's usual sparsity, and
expert disagreement; the last is a measurement rather than a self-report, which
matters given that learned confidence heads are routinely
miscalibrated~\cite{guo2017}. The composed value reaches the safety controller,
where it can veto a motor command.

The \textbf{safety controller} is deliberately not learned, holds no
parameters, and is unreachable from the plasticity controller: a policy that
rewrites its own weights at runtime must not also decide what it is permitted
to do. It enforces per-joint envelopes, named allow/clamp/veto constraints, and
a hard interlock.

\subsection{Typed sensors and actuators}

A robot does not have ``touch'' and ``motors''. Sensor specifications carry
rate, latency, noise, drift and---for chemical sensing---recovery time. The
last is not decorative: after a strong exposure a metal-oxide element is
saturated and reads high regardless of what is in front of it, so a robot
tracking a gas plume up the gradient will chase its own saturation. Fusion runs
at the slowest present stream's rate; fusing at the fastest would repeat a
two-second-old gas reading two hundred times and let it outvote the IMU.

Actuators differ more consequentially. A single uniform joint limit is wrong in
both directions on mixed hardware: the rate cap protecting a fingertip servo
renders a hydraulic hip useless, while the effort ceiling permitting the hip to
bear load will crush whatever the gripper holds. Per-type envelopes are derived
from \texttt{compliance}, \texttt{max\_effort} and bandwidth; joints that are
both stiff and strong---which will not yield to an obstacle and have the force
budget to continue through one---are flagged and held to a tighter ceiling.

\subsection{Bootstrapping from an existing teacher}

A randomly initialised architecture of this kind has no language knowledge, and
training one from scratch at useful scale is out of reach here. The repository
includes a distillation path from an open-weights teacher, with one constraint
that determines its whole design: \textbf{two models with different tokenizers
cannot be KL-matched.} BrainAI's byte tokenizer has 260 symbols against a
typical BPE vocabulary of $\sim$150k, so their output distributions live in
different spaces and the usual logit-matching objective~\cite{hinton2015} is
undefined between them.

Two modes follow. \emph{Sequence-level} distillation~\cite{kim2016} has the
teacher generate text on which the student trains with ordinary cross-entropy;
it is tokenizer-agnostic and works with any teacher, but transmits only the
sampled token rather than the distribution it came from. \emph{Logit}
distillation rebuilds the student on the teacher's vocabulary and matches full
temperature-scaled distributions, transferring considerably more at the cost of
the byte tokenizer---which for an agent also ingesting tactile and telemetry
byte streams is a substantive loss rather than a formality.

We report no distillation results. The pipeline is implemented and exercised
against a stub teacher; running it against a real one requires a GPU.

\section{Experimental Setup}

All experiments run on 8 CPU cores with no GPU. The \texttt{nano} configuration
is 5.36M parameters, 17 layers, $d_{\text{model}}=128$, $T=3$ SNN timesteps.
Pretraining uses 400 steps on Austen's \emph{Pride and Prejudice} at byte level;
the continual-learning task is Voltaire's \emph{Candide}, in French. Byte-level
English and French have clearly separated statistics, making forgetting
unambiguous. A larger multi-domain run (\S\ref{sec:breadth}) trains the same
\texttt{nano} model on 64 domains and 70.7\,MB.

We emphasise that this scale is a severe limitation, discussed in
\S\ref{sec:limitations}. It is sufficient to determine whether mechanisms
function and whether they beat their controls; it is not sufficient to
establish that the architecture is competitive at useful scale.

\section{Results}
\label{sec:results}

Table~\ref{tab:summary} summarises every controlled comparison.

\begin{table*}[t]
\centering\small
\begin{tabular}{@{}llll@{}}
\toprule
\textbf{Subsystem} & \textbf{Benchmark} & \textbf{Control} & \textbf{Result} \\
\midrule
Pathway gating & wall-clock & always-on & \win\ $2.16\times$ faster; 53.7\% of forward time \\
World model & one-step MSE & linear map / RSSM & beats RSSM, \lose\ \textbf{loses to linear} \\
World model & rollout, planning & action-cond.\ linear & ties by 3 steps; sole calibrated uncertainty \\
Curiosity & noisy-TV & pure novelty & \win\ $403\times$ preference for learnable \\
Trust grounding & broken vs.\ good sensor & fixed prior & \win\ $0.974$ vs.\ $0.104$ from identical priors \\
Reality Cortex & 5 belief scenarios & naive vote counting & \win\ wins 4 of 5 \\
Collusion defence & 600-puppet Sybil & no clustering & \win\ $1.0000$ vs.\ $0.0001$ \\
Continual learning & EN $\rightarrow$ FR forgetting & no protection & \win\ 74\% less forgetting \\
Online stability & 400 updates & no guards & \win\ $3.7\times$ less degradation \\
Planner & repeated step failure & flat planner & \win\ \texttt{remethod} vs.\ \texttt{failed} \\
Executive & deadline pressure & static priority & \win\ overtakes at the correct tick \\
Memory index & recall@4 & exact search & 94--100\% recall, \lose\ but slower \\
\textbf{AMSA} & perplexity, speed, continual & SwiGLU & \lose\ \textbf{loses on loss} \\
AMSA scale sweep & width 96 $\rightarrow$ 256 & --- & \lose\ loss gap \textbf{widens} \\
\textbf{AMSA} & sparsity / energy, 3 seeds & SwiGLU & \win\ \textbf{sparser 3/3, $\sim$14-pt lower FFN density} \\
\bottomrule
\end{tabular}
\caption{Every controlled comparison in this work. Two rows are negative and
both are load-bearing.}
\label{tab:summary}
\end{table*}

\subsection{Computational gating}

Skipping the conscious pathway is a wall-clock claim and receives a wall-clock
measurement (batch 8, sequence 128, CPU): $0.2446$\,s per forward with all rows
deliberating versus $0.1133$\,s with none, a $2.16\times$ speedup. The conscious
stack accounts for $53.7\%$ of forward time; it is the widest section and the
only expert-routed one, making it the largest single lever on inference cost.

\subsection{Continual learning}
\label{sec:continual}

\begin{table}[h]
\centering\small
\begin{tabular}{@{}lrrrr@{}}
\toprule
Arm & EN forgot & FR learned & both & vs.\ naive \\
\midrule
naive & $+0.4675$ & $+0.7591$ & $2.6882$ & --- \\
ewc   & $+0.1679$ & $+0.4543$ & $2.6908$ & 64\% less \\
full  & $\mathbf{+0.1215}$ & $+0.4380$ & $\mathbf{2.6758}$ & \textbf{74\% less} \\
\bottomrule
\end{tabular}
\caption{Forgetting in nats. \emph{both} is mean final loss over the two tasks.}
\label{tab:continual}
\end{table}

The \emph{both} column is the one to read before celebrating the forgetting
column. Protecting old weights also slows acquisition of the new task: French
improves by $0.76$ nats under \texttt{naive} but only $0.44$ under
\texttt{full}. The methods buy retention by spending plasticity at close to
$1{:}1$. The \texttt{full} arm does edge ahead on \emph{both}, by $0.012$
nats---a margin small enough on a single seed that it should be read as ``no
worse overall while forgetting 74\% less'', not as a free lunch. A forgetting
number reported without its stability--plasticity counterpart is close to
meaningless.

\subsection{Inference-time learning}

Over 400 inference-time weight updates against a control with plasticity
gating, trust region and rollback all disabled: the guarded arm degrades on the
old task by $+0.0170$ nats with parameter drift $0.2119$ and a $3\%$ rollback
rate; the control degrades by $+0.0624$ with drift $0.3151$. The guards deliver
$3.7\times$ less degradation.

We are careful about what this does \emph{not} show. The control degrades
faster; it does not diverge. At this learning rate and horizon the trust region
buys a slower slide rather than preventing a collapse. Spike rate holds at
$0.028$ in both arms, which matters independently: a spiking network that
quietly becomes dense under online updates has discarded its reason for
existing.

\subsection{Epistemics}

Against \texttt{NaiveCounter} (confidence $=$ supports$/$total):

\begin{table}[h]
\centering\small
\begin{tabular}{@{}lrr@{}}
\toprule
Scenario & Reality Cortex & Naive \\
\midrule
100 repeats vs.\ 2 good sources & $\mathbf{0.9986}$ \win & fails \\
majority is correct & $0.9857$ \win & \win \\
200 identities, one vote each & $0.0145$ \lose & fails \\
same, vs.\ protected axiom & $\mathbf{0.9956}$ \win & fails \\
200 identities with a record & $\mathbf{0.8636}$ \win & fails \\
\bottomrule
\end{tabular}
\caption{Belief quality. The two Sybil rows are the informative pair.}
\end{table}

A puppet voting exactly once is indistinguishable from a genuine one-time
witness---correctly so, as a single data point cannot support an accusation of
collusion---and gets through unless the target is a protected axiom. Once a
ring participates in several claims its shared pattern is detectable, the ring
collapses to one vote, and a \emph{single} trusted source then outweighs all
200.

Grounding removes the oracle entirely. Two sensors, both assigned prior $0.8$
by the operator, one honest and one broken, with nothing labelling either:
after 40 rounds of predictions resolving against observations, trust separates
to $0.974$ and $0.104$. A fixed prior holds both at $0.800$ indefinitely.

One methodological note: sockpuppets are assigned prior $0.62$, not $0.2$. Below
$0.5$ a source's \emph{support} is already evidence against---$\log(t/(1-t))$
is negative---so a flood of visibly untrusted accounts defeats itself and tests
nothing. Our first version of this benchmark made exactly that error and
reported a pass.

\subsection{Multi-domain breadth}
\label{sec:breadth}

The single-book runs establish that the mechanism trains; they say nothing about
breadth. We therefore trained the same \texttt{nano} model across \textbf{64
domains totalling 70.7\,MB}: code in seven languages, the natural and social
sciences, five non-Latin languages (Urdu, Hindi, Marathi, Konkani, Arabic),
Morse, binary, bytecode, assembly, geography, and an author/model knowledge
corpus. Domains are sampled at temperature $0.5$ and evaluated separately, since
one averaged number cannot separate uniform mediocrity from a model that learned
code and forgot chemistry.

Held-out loss fell from $5.49$ to \textbf{2.25 bits/byte} (mean). Selected
domains: mathematics $0.58$, IT/CS reference $0.93$, Marathi $1.31$, Hindi
$1.54$, Konkani $1.55$, Urdu $1.73$, Arabic $1.81$, Python/C++ $2.24$, and the
tiny Morse and assembly corpora at $3.0$ and $3.8$. The five non-Latin languages
were learned from scratch at the byte level with no vocabulary change---a direct
consequence of byte tokenisation. A $20\times$-boosted, EWC-protected fine-tune
(a continuation, not a retrain) then moved an author-knowledge corpus from $3.32$
to \textbf{2.59 bits/byte} while the 64-domain mean held ($1.539 \rightarrow
1.516$), demonstrating targeted acquisition without measurable forgetting.

\textbf{The ceiling is size, not data or method.} Despite falling loss
everywhere, generation remains character-level fragments: a prompt for the author
elicits ``Zahisal''---a mangled ``Zahid''---rather than a sentence. At
$1$--$2$ bits/byte the 5.36M model has learned each domain's statistics but lacks
the capacity to compose them. This is the expected regime for a model of this
size, and it is what the \texttt{compact} (104M) and \texttt{base} ($\sim$3.8B)
presets exist to move past: corpus, method and per-domain evaluation transfer
unchanged; only the parameter count and required compute increase. We report the
small-model result because it is the one that runs to completion on a CPU, and
because it isolates \emph{breadth of acquisition} from \emph{fluency}, which
scale supplies separately.

\section{Negative Results}
\label{sec:negative}

\subsection{AMSA}

The architecture proposes an activation, AMSA, whose threshold and gain are
driven by per-channel plasticity, confidence and memory state rather than by
learned constants alone:
\begin{equation}
f(x) = z\,\sigma(\alpha z)\,\big(1 - e^{-\beta|z|}\big)(1+\gamma m),
\quad z = x - \theta
\end{equation}
with $\theta = \theta_0 + \lambda p$ and
$\alpha = \text{softplus}(\alpha_0 + \delta c)$.

On static language modelling with matched parameter budgets, SwiGLU~\cite{shazeer2020}
reaches validation loss $2.3323$ against AMSA's $2.3466$ and is roughly $40\%$
faster. The natural objection is that this is the wrong benchmark: AMSA's state
buffers are inert in a single-task run, since plasticity and confidence only
carry information once a previous task's importance has been estimated. We
therefore ran the test AMSA was designed for.

\begin{table}[h]
\centering\small
\begin{tabular}{@{}lrrr@{}}
\toprule
FFN & EN forgotten & FR learned & vs.\ SwiGLU \\
\midrule
swiglu & $\mathbf{+0.8062}$ & $+1.1540$ & --- \\
amsa & $+0.8204$ & $+1.1476$ & $+0.0141$ worse \\
amsa+state & $+0.8224$ & $+1.1464$ & $+0.0161$ worse \\
\bottomrule
\end{tabular}
\caption{AMSA on continual learning, 3 seeds. Seed spread on forgetting is
$\pm0.005$--$0.008$; both gaps are outside it.}
\end{table}

AMSA loses on its home turf, and driving the adaptive state makes it
\emph{marginally worse} rather than neutral.

\paragraph{Testing the scale defence.} The remaining defence is that
\texttt{nano} is too small for the adaptive state to matter. This is testable:

\begin{table}[h]
\centering\small
\begin{tabular}{@{}rrrrrr@{}}
\toprule
width & layers & steps & swiglu & amsa & gap \\
\midrule
96 & 3 & 100 & $+0.6475$ & $\mathbf{+0.6206}$ & $-4.2\%$ \\
128 & 4 & 200 & $+0.7937$ & $+0.7898$ & $-0.5\%$ \\
192 & 4 & 200 & $+0.8313$ & $+0.8257$ & $-0.7\%$ \\
256 & 6 & 400 & $\mathbf{+0.9890}$ & $+1.0483$ & $+6.0\%$ \\
\bottomrule
\end{tabular}
\caption{The gap \emph{widens} with scale.}
\end{table}

The gap moves from $-4.2\%$ to $+6.0\%$: AMSA is genuinely \emph{better} at the
smallest setting and worse at the largest. There is a crossover, pointing the
opposite way from the hypothesis.

This must be read carefully in both directions. It does not establish that AMSA
fails at $10^9$ parameters: the range spans width 96 to 256, the two middle rows
lie inside seed noise, and width and training length were varied together rather
than independently. What it does is remove the \emph{support} for the low-data
explanation within the measured range, shifting the burden onto a larger
experiment rather than onto an assumption.

\textbf{Loss, however, is the wrong axis for AMSA.} The tests above optimise
perplexity; AMSA was designed for efficiency, and its $(1 - e^{-\beta|z|})$
factor is a small-activation suppressor. The relevant question is whether, at
comparable loss, it activates less of the network. We measure this in the real
spiking model on the 64-domain corpus, identical initialisation and batch stream,
over three seeds. On a representative seed the spike rate is $0.0282$ vs SwiGLU's
$0.0291$ and the feed-forward activation density $59.9\%$ vs $73.7\%$; across all
three seeds, at a loss difference within noise ($\leq 0.021$ nats), AMSA is
sparser in \textbf{3 of 3}, the $\sim$14-point density gap being structural---a
direct consequence of the activation's shape---while the spike-rate margin varies
by seed. In a spiking network spikes are the energy expended, so this is a
genuine advantage on the axis AMSA was built for, invisible to a
perplexity-only evaluation.

The complete picture is a quality--efficiency tradeoff, not a flat loss: SwiGLU
is marginally lower-loss and faster on dense hardware; AMSA reaches the same loss
while firing less, which is what neuromorphic hardware bills for. We retain AMSA
as the default on that basis. It remains one model size at a modest budget---AMSA
does not win on loss, the loss gap widens with scale, and the spike-rate margin,
though consistent in sign, is small---but the brain-efficiency claim now has
evidence behind it rather than only design intent.

\subsection{The world model}

Against an RSSM-family recurrent baseline and a linear map:

\begin{table}[h]
\centering\small
\begin{tabular}{@{}lrrrr@{}}
\toprule
System & persist & linear & RSSM & BrainAI \\
\midrule
linear rotation & $2.096$ & $\mathbf{0.006}$ & $0.244$ & $0.153$ \\
saturating & $1.905$ & $\mathbf{0.249}$ & $0.406$ & $0.354$ \\
\bottomrule
\end{tabular}
\end{table}

The learned model beats RSSM on both systems and loses to a linear map on both.
On the linear system this is expected---a linear map is optimal there, so
losing $24\times$ measures overhead rather than capability. On saturating
nonlinear dynamics the gap narrows to $1.4\times$ and constitutes a genuine
loss: a 16k-parameter linear layer predicts better than a 119k-parameter
dynamics head.

\paragraph{But one-step MSE is not the job.} Nothing in this architecture
requests a one-step point prediction. It requests multi-step imagination,
ranking of candidate plans, an estimate of how far ahead the model can be
trusted, and separation of learnable surprise from noise.

\begin{table}[h]
\centering\small
\begin{tabular}{@{}lrrl@{}}
\toprule
Job & linear & BrainAI & winner \\
\midrule
rollout @ 1 & $\mathbf{0.2355}$ & $0.2570$ & linear \\
rollout @ 3 & $0.3340$ & $0.3404$ & tie \\
rollout @ 6 & $0.3852$ & $0.3849$ & tie \\
rollout @ 12 & $0.4073$ & $0.4055$ & tie \\
plan: best of 8 & $38.3\%$ & $39.4\%$ & tie \\
uncertainty $r$ & n/a & $\mathbf{0.565}$ & BrainAI \\
surprise struct/noise & n/a & $-0.52/7.99$ & BrainAI \\
\bottomrule
\end{tabular}
\caption{4 seeds. Random planning baseline is $12.5\%$.}
\end{table}

The one-step deficit is real and separated at one step, and \emph{gone by
three}: both models degrade to the same place over a horizon and both plan
equally well, far above chance. The two \texttt{n/a} entries are structural
rather than unmeasured---a point prediction cannot report its own reliability,
and without that the planner cannot distinguish a trustworthy rollout from a
guess.

\section{Silent Failure Modes}
\label{sec:silent}

The following defects all had one property in common: the system continued to
run, produced plausible output, and gave no indication of malfunction. Each was
found only because a control was run or an invariant was asserted. We report
them because several generalise well beyond this codebase.

\begin{enumerate}\itemsep2pt

\item \textbf{EWC contributing nothing.} On the first continual run, the EWC
arm forgot $+0.5384$ against the no-EWC control's $+0.5385$---a $0\%$
improvement. The diagonal Fisher at convergence is $\sim\!10^{-8}$, so
$\lambda\sum F\,\Delta w^2$ evaluated to $10^{-4}$ against a task loss of
$2.5$. The penalty was arithmetically present and functionally absent. This is
indistinguishable from a working implementation without the control. The fix
was to stop guessing $\lambda$ and calibrate it, so that a $1\%$ relative drift
in every weight costs $\lambda/2$ nats and $\lambda$ carries units.

\item \textbf{Fisher normalised by its maximum.} The Fisher is heavy-tailed, so
$F/\max(F)$ is near zero for almost every weight and essentially nothing ends
up protected. Visible only by printing the protected fraction. A high quantile
is required.

\item \textbf{Query and key in different spaces.} Retrieval searched with the
raw episode summary while storage wrote \texttt{to\_key(summary)}. Cosine
similarity between vectors in different projections is meaningless; retrieval
returned confidently wrong neighbours. Best-match score rose from $0.121$ to
$1.000$ once corrected.

\item \textbf{Duplicate merging collapsing the store.} With a fixed threshold
of $0.97$, sleep merged 2042 of 2048 episodic entries. Learned summaries are
strongly collinear---genuinely distinct episodes score $\sim\!0.98$ against each
other while true duplicates sit at $\sim\!0.9998$---so a plausible-looking
constant destroyed the store while retrieval continued returning
\emph{something}. Replaced with a threshold derived from the store's own
similarity distribution.

\item \textbf{Rate limiting disabled on every call.} The guard compared
\texttt{\_last\_command.shape == command.shape}, but the reference is $(n,)$
while commands arrive as $(1,n)$. The comparison never held, so joint rate
limiting never executed. A live hazard, invisible in every test that did not
assert the ramp.

\item \textbf{No reference pose after a veto.} Following a safety veto the
controller retained no reference position, so the \emph{next} command skipped
rate limiting entirely and joints could slam immediately after a safety stop.

\item \textbf{Adaptation inert at $T=2$.} Two SNN timesteps leave the
adaptation decay multiplying only zeros, so a headline feature of the ALIF
neuron did nothing. Detected by observing that the parameter received exactly
zero gradient.

\item \textbf{A silent spiking readout.} The spiking-LSTM threshold initialised
at $0.5$ while $h = o\tanh(c)$ starts near zero, so the layer emitted all
zeros; its consumer therefore received no gradient at all. Detected by auditing
which parameters had zero gradient.

\item \textbf{Stale recurrent state across batch sizes.} Short-term memory
retained LSTM state across calls without invalidating it on a batch-size
change, silently broadcasting old state against new input and returning the
wrong batch size.

\item \textbf{A deadline making a goal less urgent.} Urgency computed as
$1/(1+k\cdot\text{remaining})$ sends distant deadlines toward zero, so attaching
a deadline to a goal made it \emph{less} likely to run than having none.

\item \textbf{Replanning checking only the immediate parent.} Alternative
methods are normally registered against the high-level task, not the leaf that
failed, so every recoverable failure was reported as fatal.

\item \textbf{Calibration measured off-distribution.} The error head trained at
unit variance and was scored across $0.3$--$2.3\times$ state magnitude, yielding
an \emph{anti}-correlation of $-0.615$ that looked like a broken mechanism.
Training across the tested range gives $+0.565$. The measurement was wrong, not
the model.

\item \textbf{An approximate index trading everything for nothing.} An initial
12-bit single-probe random-hyperplane LSH
configuration~\cite{indyk1998,charikar2002} achieved a $3.4\times$ speedup at
$19\%$ recall. Multi-probe~\cite{lv2007} reached 94--100\% recall but was then
\emph{slower} than exact search, since Python-side bucket assembly cannot beat
a single BLAS matmul below roughly 500k entries.

\item \textbf{A critic that crashed on long inputs.} A fixed window exceeded the
positional embedding of any model configured below it.

\end{enumerate}

We note that six of these---(1), (2), (4), (5), (7), (12)---produce output that
is entirely plausible. Three---(5), (6), (10)---are safety-relevant. None would
have been caught by a loss curve.

\section{Limitations}
\label{sec:limitations}

\paragraph{Scale.} Every experiment is at 5.36M parameters, including the
64-domain run of \S\ref{sec:breadth}. Continual learning is measured on two
tasks over 120 steps; online stability over 400 updates. As
\S\ref{sec:breadth} shows directly, this size learns each domain's statistics
but cannot compose them into fluent output---the deficit is capacity, not data
or method, and is expected to close only with the larger presets. Convincing
evidence of that requires GPU-scale training the CPU here cannot provide. This
is the dominant limitation and it is compute-bound rather than design-bound.

\paragraph{No robot.} Nothing here has been run on hardware, in simulation, or
against a real sensor. The body model has never calibrated against a physical
actuator; sensor specifications are representative of commodity parts rather
than measured. The safety controller enforces an envelope and named constraints,
which is a floor rather than a safety case: real deployment requires hazard
analysis, redundant sensing, and a hardware interlock independent of this
process.

\paragraph{Grounding covers a subset.} The prediction ledger removes the oracle
only for beliefs implying a checkable prediction. A claim about something the
robot never observes remains anchored to its operator prior---correct behaviour,
but it means trust is grounded for part of the knowledge base rather than all
of it.

\paragraph{Sybil resistance is a cost, not a guarantee.} Clustering detects
identities behaving identically. A patient adversary varying its puppets'
patterns, or one compromising genuinely independent sources, is not detected.

\paragraph{Triple extraction is out of scope.} The Reality Cortex reasons over
$(\text{subject},\text{relation},\text{object})$ triples that something else
must supply. Extracting reliable triples from raw perception is its own hard
problem and is not addressed.

\paragraph{Spiking energy is a hardware argument.} On CPU or GPU, binary
activations save no wall-clock time---dense matmuls execute regardless. The
$97\%$ sparsity figures are what neuromorphic hardware would exploit; here they
are a proxy.

\section{Conclusion}

BrainAI integrates spiking computation, a differentiability-split memory
hierarchy, evidence-weighted belief maintenance, predictive world modelling,
goal-directed executive control, hierarchical planning, end-to-end uncertainty,
and a non-learned safety layer into one architecture. Of thirteen controlled
comparisons, nine clearly favour the proposed component, one is mixed (the
world model ties on its downstream jobs while losing at one-step prediction),
and three favour the simpler alternative. All are reported.

We suggest the more transferable contribution may be methodological. A
component that does nothing is observationally identical to one that works,
until something is run without it. In this project that principle recovered an
EWC penalty contributing $10^{-4}$, a rate limiter disabled by a shape
comparison, a memory store collapsed to six percent of its contents, and an
uncertainty head whose apparent anti-correlation was an artefact of measuring
it off-distribution. None of these would have been visible in a loss curve, and
all of them existed in code that ran without error and produced plausible
output.

The remaining work is not architectural. It is scale, embodiment, and the
longer horizons under which the claims made here would either hold or fail.

\section*{Reproducibility}

All results are produced by scripts in the accompanying repository and run on
CPU. The test suite comprises 114 tests, including regression tests for every
defect in \S\ref{sec:silent}.

\begin{lstlisting}
python train.py --steps 400
python train_multidomain.py          # 64-domain breadth
python finetune_author.py            # author recall, a continuation
python evaluate_trained.py           # benchmarks on the trained weights
python continual.py --steps 120
python benchmarks/cognition_bench.py
python benchmarks/worldmodel_jobs.py
python benchmarks/amsa_continual.py --sweep
python benchmarks/belief_bench.py
python benchmarks/gating_bench.py
python benchmarks/stability_bench.py --control
python -m pytest tests -q
\end{lstlisting}

\begin{thebibliography}{99}\small

\bibitem{bellec2018} Bellec, G., Salaj, D., Subramoney, A., Legenstein, R., \&
Maass, W. (2018). Long short-term memory and learning-to-learn in networks of
spiking neurons. \emph{NeurIPS}.

\bibitem{burda2019} Burda, Y., Edwards, H., Pathak, D., Storkey, A., Darrell,
T., \& Efros, A. (2019). Large-scale study of curiosity-driven learning.
\emph{ICLR}.

\bibitem{charikar2002} Charikar, M. (2002). Similarity estimation techniques
from rounding algorithms. \emph{STOC}.

\bibitem{douceur2002} Douceur, J. R. (2002). The Sybil attack. \emph{IPTPS}.

\bibitem{erol1994} Erol, K., Hendler, J., \& Nau, D. (1994). HTN planning:
complexity and expressivity. \emph{AAAI}.

\bibitem{fedus2022} Fedus, W., Zoph, B., \& Shazeer, N. (2022). Switch
Transformers: scaling to trillion parameter models with simple and efficient
sparsity. \emph{JMLR}, 23.

\bibitem{guo2017} Guo, C., Pleiss, G., Sun, Y., \& Weinberger, K. Q. (2017). On
calibration of modern neural networks. \emph{ICML}.

\bibitem{hafner2019} Hafner, D., Lillicrap, T., Fischer, I., Villegas, R., Ha,
D., Lee, H., \& Davidson, J. (2019). Learning latent dynamics for planning from
pixels. \emph{ICML}.

\bibitem{hafner2020} Hafner, D., Lillicrap, T., Ba, J., \& Norouzi, M. (2020).
Dream to control: learning behaviors by latent imagination. \emph{ICLR}.

\bibitem{hinton2015} Hinton, G., Vinyals, O., \& Dean, J. (2015). Distilling
the knowledge in a neural network. \emph{arXiv:1503.02531}.

\bibitem{indyk1998} Indyk, P., \& Motwani, R. (1998). Approximate nearest
neighbors: towards removing the curse of dimensionality. \emph{STOC}.

\bibitem{kahneman2011} Kahneman, D. (2011). \emph{Thinking, Fast and Slow}.
Farrar, Straus and Giroux.

\bibitem{kim2016} Kim, Y., \& Rush, A. M. (2016). Sequence-level knowledge
distillation. \emph{EMNLP}.

\bibitem{kirkpatrick2017} Kirkpatrick, J., et al. (2017). Overcoming
catastrophic forgetting in neural networks. \emph{PNAS}, 114(13).

\bibitem{lv2007} Lv, Q., Josephson, W., Wang, Z., Charikar, M., \& Li, K.
(2007). Multi-probe LSH: efficient indexing for high-dimensional similarity
search. \emph{VLDB}.

\bibitem{maass1997} Maass, W. (1997). Networks of spiking neurons: the third
generation of neural network models. \emph{Neural Networks}, 10(9).

\bibitem{martins2016} Martins, A., \& Astudillo, R. (2016). From softmax to
sparsemax: a sparse model of attention and multi-label classification.
\emph{ICML}.

\bibitem{mcclelland1995} McClelland, J. L., McNaughton, B. L., \& O'Reilly, R.
C. (1995). Why there are complementary learning systems in the hippocampus and
neocortex. \emph{Psychological Review}, 102(3).

\bibitem{mccloskey1989} McCloskey, M., \& Cohen, N. J. (1989). Catastrophic
interference in connectionist networks: the sequential learning problem.
\emph{Psychology of Learning and Motivation}, 24.

\bibitem{neftci2019} Neftci, E. O., Mostafa, H., \& Zenke, F. (2019). Surrogate
gradient learning in spiking neural networks. \emph{IEEE Signal Processing
Magazine}, 36(6).

\bibitem{pathak2017} Pathak, D., Agrawal, P., Efros, A. A., \& Darrell, T.
(2017). Curiosity-driven exploration by self-supervised prediction.
\emph{ICML}.

\bibitem{rolnick2019} Rolnick, D., Ahuja, A., Schwarz, J., Lillicrap, T., \&
Wayne, G. (2019). Experience replay for continual learning. \emph{NeurIPS}.

\bibitem{shazeer2017} Shazeer, N., et al. (2017). Outrageously large neural
networks: the sparsely-gated mixture-of-experts layer. \emph{ICLR}.

\bibitem{shazeer2020} Shazeer, N. (2020). GLU variants improve Transformer.
\emph{arXiv:2002.05202}.

\bibitem{vaswani2017} Vaswani, A., et al. (2017). Attention is all you need.
\emph{NeurIPS}.

\bibitem{zenke2017} Zenke, F., Poole, B., \& Ganguli, S. (2017). Continual
learning through synaptic intelligence. \emph{ICML}.

\bibitem{zhou2023} Zhou, Z., et al. (2023). Spikformer: when spiking neural
network meets Transformer. \emph{ICLR}.

\end{thebibliography}

\end{document}
